\section{Introduction and Motivation}\label{sec:intro}

The past decade of artificial intelligence has been defined by the empirical scaling of dense neural architectures, culminating in Transformer-based foundation models that exhibit remarkable fluency across natural language, vision, code synthesis, and multimodal reasoning~\cite{VASWANI_TRANSFORMER}.
Yet as these models expand toward million-token context windows, autonomous lifelong deployment, and ubiquitous edge inference, the architectural assumptions underpinning dense real-valued deep learning are colliding with fundamental computational, physical, and representational limits.
Standard self-attention achieves extraordinary expressive power by retaining every historical token explicitly and computing dense pairwise interactions, an approach that strains silicon memory bandwidth and energy budgets at both cloud and edge scales.
Simultaneously, storing knowledge inside unstructured, polysemantic real-valued weights leaves contemporary models vulnerable to brittle compositional generalization, fluent factual hallucination, and catastrophic forgetting~\cite{ELHAGE_2022_SUPERPOSITION,KIRKPATRICK_2017_EWC}.

Holographic Neural Networks (HNNs) and Vector Symbolic Architectures (VSAs)---encompassing Holographic Reduced Representations (HRRs) and Hyperdimensional Computing (HDC)---offer a mathematically principled, physically grounded paradigm rooted in \emph{distributed phase-space superposition}~\cite{SUTHERLAND_1990_HNN,PLATE_1995_HRR,KANERVA_2009_HDC}.
Inspired by optical wave interference~\cite{GABOR_1948_HOLOGRAPHY}, holographic architectures encode structured associations into fixed-dimensional complex or high-dimensional vector traces via invertible algebraic binding and additive bundling.
Long regarded as a specialized branch of cognitive modeling or low-power edge classification, these algebraic principles have quietly resurfaced across the modern deep learning frontier---inside Complex-Valued Neural Networks (CVNNs)~\cite{BASSEY_CVNN_SURVEY,TRABELSI_2018_DCN}, Fourier token-mixing layers~\cite{LEE_THORP_FNET}, kernelized linear attention~\cite{KATHAROPOULOS_2020_LINEAR_TRANSFORMER}, Rotary Position Embeddings~\cite{SU_2024_ROPE}, and diagonal complex State-Space Models (SSMs)~\cite{GU_SSM_S4_MAMBA}.
This survey unifies these historically fragmented traditions under a single rigorous mathematical and comparative lens, evaluates their contemporary applications and hardware co-design synergies, and charts their potential for next-generation AI systems.

% ------------------------------------------------------------------------------
\subsection{The Scaling Wall: Compute, Memory, and Energy Bottlenecks in Modern AI}\label{sec:intro:bottlenecks}
% ------------------------------------------------------------------------------

Three interlocking bottlenecks expose the structural limits of prevailing dense real-valued neural architectures as workloads scale in context length, deployment autonomy, and reasoning rigor.

\paragraph{Quadratic Attention and Linearly Growing KV-Caches.}
In a canonical Transformer layer operating over a sequence of length $N$ and hidden dimension $D$, dense scaled dot-product self-attention evaluates all pairwise query--key similarities, incurring $\bigo{N^2 D}$ floating-point operations and $\bigo{N^2}$ pairwise score storage~\cite{VASWANI_TRANSFORMER}.
While hardware-aware tiled kernels mitigate the transient materialization of the $N \times N$ score matrix during training and prefill, autoregressive token generation exposes an equally severe bottleneck: to avoid recomputing past projections, inference engines must retain an explicit key--value (KV) cache whose memory footprint grows linearly as $\bigo{ND}$ per layer (specifically $\bigo{L H N d_h}$ across $L$ layers and $H$ heads of width $d_h$).
At context lengths reaching $N \ge 10^5\text{--}10^6$ tokens, the KV-cache dominates accelerator memory capacity, eclipsing even the model's static parameter footprint and forcing aggressive batch-size reduction or lossy token eviction.

\paragraph{The Memory-Bandwidth Wall and Energy Constraints.}
Over the past decade, raw peak floating-point throughput on specialized AI accelerators (GPU Tensor Cores and systolic TPUs) has scaled far more rapidly than off-chip high-bandwidth memory (HBM) bandwidth and interconnect energy efficiency.
During autoregressive decoding, generating each new token requires streaming the entire $\bigo{ND}$ KV-cache alongside static weights across off-chip memory buses to perform only $\bigo{ND}$ arithmetic operations, locking frontier inference into a strictly memory-bandwidth-bound regime in which arithmetic units sit largely idle.
The resulting escalation in energy per token ($\mathrm{J/tok}$) and data-movement overhead poses acute economic and environmental challenges in cloud data centers while rendering dense attention outright infeasible on milliwatt-budget wearable sensors, micro-drones, and always-on Internet-of-Things (IoT) edge nodes~\cite{RAHIMI_HDC_EDGE}.
Moreover, because modern hardware stacks are co-optimized almost exclusively for low-precision real-valued general matrix multiplication (GEMM), alternative representational paradigms face a structural \emph{hardware lottery} penalty unless algorithm design and physical substrate engineering evolve in tandem~\cite{HOOKER_2021_HARDWARE_LOTTERY}.

\paragraph{Representational Limits of Unstructured Dense Weights.}
Beyond systems efficiency, purely dense, real-valued representations exhibit persistent cognitive and reliability failure modes:
\begin{itemize}
    \item \emph{The binding problem and brittle compositionality:} Without an explicit algebraic operator that separates relational \emph{roles} from entity \emph{fillers}~\cite{SMOLENSKY_1990_TPR}, dense networks store features via implicit polysemantic superposition~\cite{ELHAGE_2022_SUPERPOSITION}. Co-occurring attributes become statistically entangled, impairing systematic out-of-distribution recombination and causing attribute-swapping errors.
    \item \emph{Fluent hallucination and uncalibrated confidence:} Standard autoregressive softmax decoders normalize logits over the entire vocabulary regardless of whether a query lies within the model's stored knowledge, yielding confident statistical confabulations without an intrinsic signal of memory absence or retrieval collision.
    \item \emph{Catastrophic forgetting under continual updates:} Because gradient updates perturb shared real-valued synaptic weights globally, adapting a network to a new task or updating a single factual record frequently overwrites previously acquired knowledge unless mitigated by expensive episodic replay, architectural expansion, or strong plasticity-limiting regularization~\cite{KIRKPATRICK_2017_EWC}.
\end{itemize}

% ------------------------------------------------------------------------------
\subsection{From Optical Holography to Latent Superposition}\label{sec:intro:holography}
% ------------------------------------------------------------------------------

Holographic computation addresses these computational and representational bottlenecks by replacing localized slot-based storage with \emph{distributed wave-like superposition}.
In optical holography, first formalized by Dennis Gabor~\cite{GABOR_1948_HOLOGRAPHY}, a coherent reference beam interferes with an object wavefield across an entire photosensitive medium, recording both relative phase and amplitude as a global interference fringe pattern.
Illuminating any sub-region of the recorded hologram with the reference beam reconstructs the entire stored wavefield, degrading gracefully in signal resolution rather than suffering localized data loss.
Longuet-Higgins subsequently translated this physical principle into frequency-domain associative temporal recall~\cite{LONGUET_HIGGINS_1968_HOLOPHONE}, laying the conceptual groundwork for neural models that store and query information via algebraic interference rather than explicit addressing.

Modern HNNs and VSAs lift this physical analogy into high-dimensional real ($\R^D$) or complex Hilbert spaces ($\C^D$), and most cleanly onto the unitary complex torus $\mathbb{T}^D = \{\vect{z} \in \C^D : \abs{z_k} = 1\}$~\cite{SUTHERLAND_1990_HNN,PLATE_1995_HRR,KANERVA_2009_HDC}.
In polar (phasor) coordinates, each component $x_k = \lambda_k e^{i\theta_k}$ decouples symbolic or continuous feature value (phase angle $\theta_k \in [-\pi, \pi)$) from confidence or salience (magnitude $\lambda_k \ge 0$).
Four mathematical and physical properties distinguish this holographic paradigm from conventional dense architectures:
\begin{enumerate}
    \item \emph{Distributed non-locality and fault tolerance:} By the concentration of measure in high dimensions, independent random vectors in $\mathbb{T}^D$ or $\R^D$ are mutually \emph{quasi-orthogonal} with exponentially small overlap probability, accommodating exponentially many distinct atomic symbols in $D$ dimensions. Because every stored association is distributed uniformly across all $D$ coordinates, holographic traces exhibit graceful degradation under coordinate erasures, aggressive low-bit quantization, memory bit-flips, and analog hardware drift~\cite{KANERVA_2009_HDC,KARUNARATNE_2020_PCM_HDC}.
    \item \emph{Constant-size superposition with analytical capacity scaling:} Rather than appending each new key--value pair $(\vect{k}_n, \vect{v}_n)$ to a growing $\bigo{ND}$ cache, holographic bundling enfolds $N$ bound pairs into a single composite trace $\vect{m} = \sum_{n=1}^{N} (\vect{k}_n \cconv \vect{v}_n) \in \C^D$ of strictly $\bigo{D}$ size---that is, $\bigo{1}$ memory with respect to $N$. This compression trades exact non-parametric storage for zero-mean additive crosstalk noise $\vect{\eta}$ whose magnitude is governed by the exact Signal-to-Noise Ratio law $\SNR \propto D/N$~\cite{PLATE_1995_HRR,FRADY_CAPACITY_VSA}. Coupled with a clean-up codebook of size $M$, reliable discrete retrieval scales as $N_{\max} = \Theta(D / \ln M)$, exhibiting only logarithmic sensitivity to vocabulary size.
    \item \emph{Native dimension-preserving compositionality:} Binding via circular convolution ($\cconv$) in the time domain---or equivalently $\bigo{D}$ elementwise Hadamard phase addition ($\hadamard$) in the Fourier domain---maps $\C^D \times \C^D \to \C^D$ without inflating dimensionality, resolving the exponential tensor-rank explosion of Tensor Product Representations~\cite{SMOLENSKY_1990_TPR,PLATE_1995_HRR}. Querying the composite trace via the Hilbert-space adjoint, circular correlation ($\ccorr$), algebraically unbinds specific role--filler associations in $\bigo{D}$--$\bigo{D \log D}$ time while enabling one-shot online insertion and deletion without gradient backpropagation.
    \item \emph{Direct physical isomorphism with wave and analog substrates:} Unlike real-valued GEMM-centric networks, complex-phasor holographic algebra maps directly onto passive optical physics: wavefront interference executes bundling, programmable phase masks execute frequency-domain binding $\hadamard$, and thin spherical lenses compute exact spatial Fourier transforms $\DFT / \IDFT$ at the speed of light~\cite{SHEN_PHOTONIC_MZI,LIN_DIFFRACTIVE_OPTICAL_NN,WETZSTEIN_2020_OPTICAL_AI}. Conversely, discrete bipolar and binary VSA specializations collapse binding to bitwise XOR and similarity to hardware popcount, unlocking sub-microjoule edge inference and resilient analog in-memory computing~\cite{RAHIMI_HDC_EDGE,KARUNARATNE_2020_PCM_HDC}.
\end{enumerate}

% ------------------------------------------------------------------------------
\subsection{Survey Methodology and Scope}\label{sec:intro:methodology}
% ------------------------------------------------------------------------------

\paragraph{Literature Search Strategy and Selection Criteria.}
To assemble a representative, reproducible synthesis across neural, symbolic, and photonic communities, we conducted structured literature searches across four primary academic databases---\emph{IEEE Xplore}, the \emph{ACM Digital Library (ACM DL)}, \emph{arXiv}, and \emph{Google Scholar}---supplemented by backward and forward citation chaining from seminal works~\cite{GABOR_1948_HOLOGRAPHY,SUTHERLAND_1990_HNN,PLATE_1995_HRR,KANERVA_2009_HDC,RAMSAUER_MODERN_HOPFIELD,GU_SSM_S4_MAMBA}, with a literature cutoff date of \textbf{mid-2026}.
Search queries combined paradigm terms (\emph{holographic neural network}, \emph{vector symbolic architecture}, \emph{holographic reduced representation}, \emph{hyperdimensional computing}, \emph{complex-valued neural network}) with architectural and substrate terms (\emph{circular convolution}, \emph{state-space model}, \emph{linear attention}, \emph{Hopfield}, \emph{KV-cache}, \emph{photonic/diffractive neural network}).
We included peer-reviewed journal and conference papers as well as widely cited archival preprints that contribute (i)~formal algebraic or statistical capacity results, (ii)~differentiable sequence or memory architectures with explicit superposition structure, (iii)~validated neuro-symbolic, edge, or optical/analog hardware realizations, or (iv)~standardized long-context recall benchmarks.
We excluded purely optical imaging studies lacking neural or associative computation, incremental application-specific HDC classifiers without novel algebraic or systems insights, and non-archival blog posts.

\noindent\textit{Methodological limitation.} This is a structured, narrative interdisciplinary survey, not a PRISMA-style systematic review. Per-query retrieval dates, hit counts, deduplication records, and screening decisions were not retained in a machine-readable search log; consequently, exact reconstruction of the candidate-paper pool cannot be guaranteed. The cited literature should not be interpreted as an exhaustive census of all relevant work.

\paragraph{Taxonomy of Cross-Paradigm Links.}
Because interdisciplinary surveys risk conflating exact mathematics with loose metaphor, we distinguish three levels of analytical rigor throughout Sections~\ref{sec:history}--\ref{sec:potentials}:
\begin{itemize}
    \item \emph{Exact algebraic equivalence:} A formal identity holding up to isomorphism, basis change, or parameter restriction---for instance, circular convolution/correlation versus spectral Hadamard modulation via the DFT (Proposition~\ref{prop:conv_thm}, Lemma~\ref{lem:ccorr_props}), single-step Modern Hopfield updates versus scaled dot-product attention~\eqref{eq:modern_hopfield_update}--\eqref{eq:transformer_attn}, and causal linear attention's outer-product state~\eqref{eq:comp_linear_attn_state} as a real-valued specialization of Sutherland's correlation matrix~\eqref{eq:hnn_corr_matrix}.
    \item \emph{Structural analogy:} A shared governing mechanism across distinct domains, parameterizations, or learning rules---for example, diagonal complex SSM phase rotations $e^{i\omega_d(t-\tau)}$ acting as a continuous-frequency analog of VSA permutation encoding, optical wave interference inspiring complex phasor bundling, and DeltaNet/Gated DeltaNet rank-one erasure paralleling Sutherland's iterative delta rule.
    \item \emph{Architectural correspondence:} A functional or systems-level pairing occupying an analogous role in a larger compute graph without claiming algebraic identity---such as FNet's unparameterized DFT token mixing relative to spectral holographic transforms, or local-plus-recurrent hybrids (Griffin, Titans) relative to the hybrid holographic KV-cache blueprint in Section~\ref{sec:potentials:kvcache}.
\end{itemize}

\paragraph{Established Literature Results vs.\ Novel Syntheses.}
To keep provenance transparent, we explicitly separate established results from novel syntheses below:
\emph{Established literature results} (unified here under a single complex-phasor notation) include quasi-orthogonality on $\mathbb{T}^D$ (Proposition~\ref{prop:quasi_ortho}), DFT convolution--correlation duality (Proposition~\ref{prop:conv_thm}, Lemma~\ref{lem:ccorr_props}), the $\SNR = \Theta(D/N)$ scaling law and asymptotic clean-up capacity bounds (Theorem~\ref{thm:snr_capacity} and Corollary~\ref{cor:cleanup_capacity}, unifying Plate~\cite{PLATE_1995_HRR} and Frady et al.~\cite{FRADY_CAPACITY_VSA}), and the Modern Hopfield--Transformer equivalence~\cite{RAMSAUER_MODERN_HOPFIELD}.
\emph{Novel syntheses and conceptual blueprints} introduced in this survey comprise the side-by-side exact finite-$D$ FHRR vs.\ real Gaussian HRR discrimination-variance comparison with self-contained Monte Carlo verification (Theorem~\ref{thm:snr_capacity} synthesis, Figure~\ref{fig:capacity}, Appendix~\ref{app:simulation}), the three-regime parametric versus non-parametric duality framing (Proposition~\ref{prop:superposition_duality}), the spectral-torus hybrid $\bigo{1}$ KV-cache blueprint (Section~\ref{sec:potentials:kvcache}, Figure~\ref{fig:hybrid_arch}), the retrieval-coherence epistemic diagnostic $\rho_{\mathrm{coh}}$ (Section~\ref{sec:potentials:reasoning}), and the unitary spectral-torus formulation of continual parameter superposition (Section~\ref{sec:potentials:continual}).

% ------------------------------------------------------------------------------
\subsection{Positioning, Contributions, and Organization}\label{sec:intro:contributions}
% ------------------------------------------------------------------------------

\paragraph{Positioning Relative to Prior Surveys.}
Several valuable surveys have examined constituent branches of this landscape in isolation.
Schlegel et al.~\cite{SCHLEGEL_VSA_COMPARISON} and Kleyko et al.~\cite{KLEYKO_VSA_SURVEY} provide authoritative taxonomies and algebraic comparisons of Vector Symbolic Architectures and Hyperdimensional Computing, while Ge and Parhi~\cite{GE_2020_HDC_CLASSIFICATION} focus specifically on HDC classification pipelines.
In parallel, Bassey et al.~\cite{BASSEY_CVNN_SURVEY} survey deep Complex-Valued Neural Networks (CVNNs) and Wirtinger optimization, whereas Wetzstein et al.~\cite{WETZSTEIN_2020_OPTICAL_AI} and McMahon~\cite{MCMAHON_2023_PHYSICS_OPTICAL_NN} review the physics and engineering of optical and photonic neural accelerators.
What remains missing in the literature is a unified, mathematically cohesive synthesis that connects all four domains---classical 1990s Holographic Neural Networks, algebraic VSAs, contemporary deep spectral and state-space sequence models, and photonic/analog hardware substrates---and evaluates them directly against the memory, reasoning, and continual-learning bottlenecks of modern foundation models.

\paragraph{Scope and Organizing Taxonomy.}
This survey establishes two complementary taxonomic frameworks that structure the entire manuscript:
\begin{itemize}
    \item \emph{A three-wave genealogical taxonomy} (Section~\ref{sec:history}, Figure~\ref{fig:genealogy}) tracing the intellectual convergence from \emph{First-Wave} optical analogies and Sutherland's complex-correlation HNNs~\cite{SUTHERLAND_1990_HNN}, through \emph{Second-Wave} algebraic VSAs (Plate's HRR/FHRR~\cite{PLATE_1995_HRR} and Kanerva's HDC~\cite{KANERVA_2009_HDC}), to \emph{Third-Wave} end-to-end differentiable CVNNs~\cite{TRABELSI_2018_DCN}, unitary recurrences~\cite{ARJOVSKY_2016_URNN}, Fourier token mixers~\cite{LEE_THORP_FNET}, linear attention~\cite{KATHAROPOULOS_2020_LINEAR_TRANSFORMER}, and structured/selective State-Space Models~\cite{GU_SSM_S4_MAMBA}.
    \item \emph{A cross-paradigm memory and hardware taxonomy} (Sections~\ref{sec:comparison}--\ref{sec:applications}) contrasting \emph{parametric holographic superposition} against \emph{non-parametric explicit attention} (Modern Hopfield Networks~\cite{RAMSAUER_MODERN_HOPFIELD} and softmax Transformers~\cite{VASWANI_TRANSFORMER}) along six rigorous axes---memory footprint, retrieval compute, training plasticity, symbolic compositionality, capacity/noise degradation, and hardware affinity (Table~\ref{tab:paradigm_comparison})---paired with a three-substrate co-design analysis across digital edge CMOS, analog in-memory crossbars, and coherent photonics (Table~\ref{tab:hardware_substrates}).
\end{itemize}
Throughout, we intentionally treat holographic superposition not as a panacea that replaces dense attention unconditionally, but as a precisely characterized mathematical regime whose Shannon-capacity boundaries dictate where pure holographic systems excel and where hybrid holographic--attention architectures are required.

\paragraph{Summary of Contributions.}
Concretely, this survey makes five main contributions:
\begin{enumerate}
    \item \textbf{Unified Complex-Phasor Mathematical Framework (Section~\ref{sec:math}):} We develop a self-contained formulation over $\C^D$ and $\mathbb{T}^D$ that subsumes Sutherland HNNs, FHRR, real HRR, MAP, and BSC; proves exact time--frequency duality between $\cconv/\ccorr$ and $\hadamard$ (Proposition~\ref{prop:conv_thm}, Lemma~\ref{lem:ccorr_props}); derives exact crosstalk statistics clarifying the conditional factor-of-two discrimination $\SNR$ advantage of unitary complex phasors over real Gaussian HRR (Theorem~\ref{thm:snr_capacity}) and asymptotic clean-up retrieval capacity (Corollary~\ref{cor:cleanup_capacity}); and validates these laws via reproducible Monte Carlo simulation (Figures~\ref{fig:algebra}--\ref{fig:capacity}, Appendix~\ref{app:simulation}).
    \item \textbf{Genealogical Synthesis Across Four Decades (Section~\ref{sec:history}):} We trace the three historical waves linking optical holography and 1990s HNNs to modern CVNNs, FNet, linear attention, and SSMs, systematically separating exact algebraic equivalences (e.g., Sutherland's correlation matrix and causal linear attention) from structural analogies (e.g., complex-diagonal state-space recurrences and delta-rule erasure) and architectural correspondences (Figure~\ref{fig:genealogy}).
    \item \textbf{Multi-Axis Comparative Evaluation and Capacity Duality (Section~\ref{sec:comparison}):} We benchmark HNNs, HRR/HDC, Classical/Modern Hopfield Networks, and Transformer attention across six analytical axes (Table~\ref{tab:paradigm_comparison}), formalize the information-theoretic duality between parametric superposition and non-parametric retention (Proposition~\ref{prop:superposition_duality}), and delineate the operational regimes in which each paradigm dominates.
    \item \textbf{Applications and Algorithm--Hardware Co-Design Synthesis (Section~\ref{sec:applications}):} We survey deployed holographic systems across neuro-symbolic reasoning (knowledge-graph HolE~\cite{NICKEL_2016_HOLE}, resonator factorization~\cite{FRADY_2020_RESONATOR,KENT_2020_RESONATOR}, and NVSA~\cite{HERSCHE_2023_NVSA}), ultra-low-power edge biosignal, robotics, and genomics accelerators~\cite{RAHIMI_HDC_EDGE,NEUBERT_2019_HDC_ROBOTICS,KIM_2020_GENIEHD,KARUNARATNE_2020_PCM_HDC}, and coherent photonic architectures including Mach--Zehnder interferometer meshes~\cite{SHEN_PHOTONIC_MZI,CLEMENTS_2016_MZI} and diffractive optical networks~\cite{LIN_DIFFRACTIVE_OPTICAL_NN} (Table~\ref{tab:hardware_substrates}).
    \item \textbf{Forward-Looking Research Roadmap for Foundation-Scale AI (Section~\ref{sec:potentials}):} We formulate architectural blueprints for hierarchical holographic $\bigo{1}$ KV-caches integrated with local sliding-window attention (Figure~\ref{fig:hybrid_arch}), crosstalk-aware epistemic uncertainty metrics ($\rho_{\mathrm{coh}}$) and auditable external indices for hallucination-reduced reasoning, and orthogonal phase-subspace parameter superposition for continual learning, while identifying the core information-theoretic, optimization, systems, and benchmarking challenges ahead.
\end{enumerate}

\paragraph{Paper Organization.}
The remainder of this survey is organized as follows.
Section~\ref{sec:math} develops the mathematical foundations of complex phasor spaces, algebraic binding/unbinding, crosstalk noise, and capacity bounds.
Section~\ref{sec:history} charts the three-wave historical evolution and architectural taxonomy from optical holography to modern spectral and state-space models.
Section~\ref{sec:comparison} conducts a systematic cross-paradigm comparison against Hopfield associative memories and Transformer self-attention.
Section~\ref{sec:applications} examines current deployments in neuro-symbolic AI, low-power edge and in-memory computing, and coherent photonic accelerators.
Section~\ref{sec:potentials} explores future potentials for foundation-scale models---including holographic KV-caches, verifiable reasoning, and continual learning---alongside open scaling bottlenecks.
Finally, Section~\ref{sec:conclusion} synthesizes our mathematical, architectural, and hardware co-design findings and outlines strategic research priorities.
